
\documentclass{article}
\usepackage{colortbl}
\usepackage{makecell}
\usepackage{multirow}
\usepackage{supertabular}

\begin{document}

\newcounter{utterance}

\centering \large Interaction Transcript for game `cladder', experiment `full\_v1.5\_default', episode 1143 with Qwen3.6{-}27B.
\vspace{24pt}

{ \footnotesize  \setcounter{utterance}{1}
\setlength{\tabcolsep}{0pt}
\begin{supertabular}{c@{$\;$}|p{.15\linewidth}@{}p{.15\linewidth}p{.15\linewidth}p{.15\linewidth}p{.15\linewidth}p{.15\linewidth}}
    \# & \multicolumn{2}{c}{Player} && \multicolumn{2}{c}{Game Master} \\
    \hline

    \theutterance \stepcounter{utterance}  
    & & & \multicolumn{4}{p{0.6\linewidth}}{
        \cellcolor[rgb]{0.9,0.9,0.9}{
            \makecell[{{p{\linewidth}}}]{
                \texttt{\tiny{[P1$\langle$GM]}}
                \texttt{You are an expert in causal inference. The following question is not a typical commonsense query, but rather a meticulously designed question created by a professor specializing in causal inference, intended to assess the students' mastery of the course content.} \\
\texttt{Imagine a self{-}contained, hypothetical world with only the following conditions, and without any unmentioned factors or causal relationships: Kidney stone size has a direct effect on treatment and recovery. Treatment has a direct effect on recovery. Method 1: We look directly at how treatment correlates with recovery in general. Method 2: We look at this correlation case by case according to kidney stone size. To understand how treatment affects recovery, is it more correct to use the Method 1 than Method 2?} \\
\texttt{Start your answer with 'Yes' or 'No', followed by additional reasoning or evidence to support your explanation.} \\
            }
        }
    }
     \\ \\

    \theutterance \stepcounter{utterance}  
    & \multicolumn{4}{p{0.6\linewidth}}{
        \cellcolor[rgb]{0.561,0.737,0.561}{
            \makecell[{{p{\linewidth}}}]{
                \texttt{\tiny{[P1$\rangle$GM]}}
                \texttt{No} \\
\\ 
\texttt{**Reasoning:**} \\
\\ 
\texttt{This scenario is a classic illustration of **Simpson’s Paradox** in causal inference. To determine the correct method, we must analyze the causal structure provided:} \\
\\ 
\texttt{1.  **Causal Structure:**} \\
\texttt{    *   Kidney Stone Size (\$S\$) \$\textbackslash{}rightarrow\$ Treatment (\$T\$)} \\
\texttt{    *   Kidney Stone Size (\$S\$) \$\textbackslash{}rightarrow\$ Recovery (\$R\$)} \\
\texttt{    *   Treatment (\$T\$) \$\textbackslash{}rightarrow\$ Recovery (\$R\$)} \\
\\ 
\texttt{    In this structure, Kidney Stone Size (\$S\$) is a **common cause** (or confounder) of both Treatment (\$T\$) and Recovery (\$R\$). It is not a mediator (because \$T\$ does not cause \$S\$) and it is not a collider.} \\
\\ 
\texttt{2.  **Analysis of Method 1 (Marginal Association):**} \\
\texttt{    Method 1 looks at the direct correlation between Treatment and Recovery without adjusting for Kidney Stone Size. Because \$S\$ influences both \$T\$ and \$R\$, the observed association between \$T\$ and \$R\$ will be biased by the confounding effect of \$S\$. For example, if larger stones lead to a different treatment choice and also independently lead to worse recovery, the marginal correlation might suggest a treatment is harmful or beneficial when, in reality, the effect is driven by the stone size. This is known as **confounding bias**.} \\
\\ 
\texttt{3.  **} \\
            }
        }
    }
    & & \\ \\

    \theutterance \stepcounter{utterance}  
    & & & \multicolumn{2}{p{0.3\linewidth}}{
        \cellcolor[rgb]{0.9,0.9,0.9}{
            \makecell[{{p{\linewidth}}}]{
                \texttt{\tiny{[GM$|$GM]}}
                \texttt{No} \\
            }
        }
    }
    & & \\ \\

    \theutterance \stepcounter{utterance}  
    & & & \multicolumn{2}{p{0.3\linewidth}}{
        \cellcolor[rgb]{0.9,0.9,0.9}{
            \makecell[{{p{\linewidth}}}]{
                \texttt{\tiny{[GM$|$GM]}}
                \texttt{no} \\
            }
        }
    }
    & & \\ \\

    \theutterance \stepcounter{utterance}  
    & & & \multicolumn{2}{p{0.3\linewidth}}{
        \cellcolor[rgb]{0.9,0.9,0.9}{
            \makecell[{{p{\linewidth}}}]{
                \texttt{\tiny{[GM$|$GM]}}
                \texttt{game\_result = WIN} \\
            }
        }
    }
    & & \\ \\

\end{supertabular}
}

\end{document}
